\documentclass[10pt]{article}
\usepackage[margin=0.6in]{geometry}
\usepackage{titlesec}
\titlespacing*{\section}{0pt}{6pt}{3pt}
\titlespacing*{\subsection}{0pt}{4pt}{1pt}
\titleformat*{\section}{\large\bfseries}
\titleformat*{\subsection}{\normalsize\bfseries}
\usepackage{amsmath, amssymb}
\usepackage{xcolor}
\usepackage{tcolorbox}

\newcommand{\pd}[2]{\frac{\partial #1}{\partial #2}}
\newcommand{\pdd}[2]{\frac{\partial^2 #1}{\partial #2^2}}
\newcommand{\dd}[2]{\frac{d #1}{d #2}}

\title{Projectile Motion\\[3pt]
\large From Newton's second law to range, height and time of flight}
\author{Derivation Notes --- Claude Code session \texttt{meetingproject}, via RoboMeet}
\date{September 2026}

\begin{document}
\begin{center}{\Large\bfseries Projectile Motion}\\[2pt]{\normalsize From Newton's second law to range, height and time of flight}\\[2pt]{\small Derivation notes --- Claude Code session \texttt{meetingproject}, via RoboMeet --- September 2026}\end{center}
\vspace{-0.6em}
\setlength{\parskip}{1pt}
\setlength{\abovedisplayskip}{3pt}\setlength{\belowdisplayskip}{3pt}
\small

%% ============================================================
\section{Setup}

A point mass $m$ is launched from the origin with speed $v_0$ at angle $\theta$ above the horizontal.
Gravity $g$ acts downward; air resistance is neglected.

\begin{tcolorbox}[colback=blue!5, colframe=blue!50!black, title={Governing equations (Newton's second law, per component)}, boxsep=2pt, top=2pt, bottom=2pt]
\begin{align}
  \textbf{Horizontal:}& \quad m\,\ddot{x} = 0 \label{eq:x}\\[2pt]
  \textbf{Vertical:}& \quad m\,\ddot{y} = -m g \label{eq:y}\\[2pt]
  \textbf{Initial conditions:}& \quad x(0)=0,\; y(0)=0,\quad
    \dot{x}(0)=v_0\cos\theta,\; \dot{y}(0)=v_0\sin\theta
\end{align}
\end{tcolorbox}

\noindent
The two directions are \emph{uncoupled}: gravity touches only $y$. Each equation integrates on its own.

%% ============================================================
\section{Integrate twice}

\begin{align}
  \dot{x}(t) &= v_0\cos\theta, &
  x(t) &= v_0\cos\theta\; t \label{eq:xt}\\[3pt]
  \dot{y}(t) &= v_0\sin\theta - g t, &
  y(t) &= v_0\sin\theta\; t - \tfrac{1}{2} g t^2 \label{eq:yt}
\end{align}

\subsection*{Trajectory: eliminate $t$}
From \eqref{eq:xt}, $t = x/(v_0\cos\theta)$. Substitute into \eqref{eq:yt}:
\begin{equation}
  \boxed{\,y(x) = x\tan\theta - \frac{g}{2 v_0^2\cos^2\theta}\,x^2\,}
  \qquad\text{a parabola through the origin.}
  \label{eq:traj}
\end{equation}

%% ============================================================
\section{Three results}

\subsection*{Time of flight: $y = 0$ again}
\begin{equation}
  v_0\sin\theta\; T - \tfrac{1}{2} g T^2 = 0
  \;\Longrightarrow\;
  \boxed{\,T = \frac{2 v_0\sin\theta}{g}\,}
  \label{eq:T}
\end{equation}

\subsection*{Maximum height: $\dot{y} = 0$}
The peak is reached at $t^\ast = v_0\sin\theta/g = T/2$ (by symmetry). Then
\begin{equation}
  H = v_0\sin\theta\, t^\ast - \tfrac{1}{2} g\, t^{\ast 2}
  \;\Longrightarrow\;
  \boxed{\,H = \frac{v_0^2\sin^2\theta}{2g}\,}
  \label{eq:H}
\end{equation}

\subsection*{Range: $x$ at $t = T$}
\begin{equation}
  R = v_0\cos\theta\cdot\frac{2 v_0\sin\theta}{g}
  \;\Longrightarrow\;
  \boxed{\,R = \frac{v_0^2\sin 2\theta}{g}\,}
  \qquad (\sin 2\theta = 2\sin\theta\cos\theta)
  \label{eq:R}
\end{equation}

%% ============================================================
\section{Observations}

\begin{tcolorbox}[colback=green!5, colframe=green!40!black, title={What the formulas say}, boxsep=2pt, top=2pt, bottom=2pt]
\begin{itemize}
  \setlength{\itemsep}{2pt}
  \item \textbf{Best angle.} $R$ is largest when $\sin 2\theta = 1$, i.e.\ $\theta = 45^\circ$, giving $R_{\max} = v_0^2/g$.
  \item \textbf{Complementary angles.} $\theta$ and $90^\circ-\theta$ share the same range ($\sin 2\theta$ is symmetric about $45^\circ$) but not the same height or flight time.
  \item \textbf{Scaling.} Doubling $v_0$ quadruples both $R$ and $H$; doubling $g$ halves them. Mass never appears.
  \item \textbf{Symmetry.} Rise and fall take equal time $T/2$; the speed at landing equals the launch speed.
  \item \textbf{Check.} $\theta = 90^\circ$: $R = 0$, $H = v_0^2/2g$, $T = 2v_0/g$ (straight up and down). $\theta = 0$: $T = 0$, nothing leaves the ground.
\end{itemize}
\end{tcolorbox}

\end{document}
